
\documentclass{article}
\usepackage{colortbl}
\usepackage{makecell}
\usepackage{multirow}
\usepackage{supertabular}

\begin{document}

\newcounter{utterance}

\centering \large Interaction Transcript for game `cladder', experiment `full\_v1.5\_default', episode 9 with llama{-}3.1{-}8b{-}instruct.
\vspace{24pt}

{ \footnotesize  \setcounter{utterance}{1}
\setlength{\tabcolsep}{0pt}
\begin{supertabular}{c@{$\;$}|p{.15\linewidth}@{}p{.15\linewidth}p{.15\linewidth}p{.15\linewidth}p{.15\linewidth}p{.15\linewidth}}
    \# & \multicolumn{2}{c}{Player} && \multicolumn{2}{c}{Game Master} \\
    \hline

    \theutterance \stepcounter{utterance}  
    & & & \multicolumn{4}{p{0.6\linewidth}}{
        \cellcolor[rgb]{0.9,0.9,0.9}{
            \makecell[{{p{\linewidth}}}]{
                \texttt{\tiny{[P1$\langle$GM]}}
                \texttt{You are an expert in causal inference. The following question is not a typical commonsense query, but rather a meticulously designed question created by a professor specializing in causal inference, intended to assess the students' mastery of the course content.} \\
\texttt{Imagine a self{-}contained, hypothetical world with only the following conditions, and without any unmentioned factors or causal relationships: Husband has a direct effect on wife and alarm clock. Wife has a direct effect on alarm clock. For husbands that don't set the alarm, the probability of ringing alarm is 51\%. For husbands that set the alarm, the probability of ringing alarm is 49\%. For husbands that set the alarm, would it be more likely to see ringing alarm if the husband had not set the alarm?} \\
\texttt{Start your answer with 'Yes' or 'No', followed by additional reasoning or evidence to support your explanation.} \\
            }
        }
    }
     \\ \\

    \theutterance \stepcounter{utterance}  
    & \multicolumn{4}{p{0.6\linewidth}}{
        \cellcolor[rgb]{0.561,0.737,0.561}{
            \makecell[{{p{\linewidth}}}]{
                \texttt{\tiny{[P1$\rangle$GM]}}
                \texttt{No.} \\
\\ 
\texttt{To answer this question, we need to consider the causal relationships and conditional probabilities given in the problem. } \\
\\ 
\texttt{Let's denote the event of the husband setting the alarm as 'A' and the event of the alarm ringing as 'R'. We are given the following conditional probabilities:} \\
\\ 
\texttt{{-} P(R|A) = 0.49 (probability of the alarm ringing given that the husband set the alarm)} \\
\texttt{{-} P(R|\textasciitilde{}A) = 0.51 (probability of the alarm ringing given that the husband did not set the alarm)} \\
\\ 
\texttt{We are asked to compare the probability of the alarm ringing if the husband had not set the alarm (P(R|\textasciitilde{}A)) with the probability of the alarm ringing if the husband had set the alarm (P(R|A)).} \\
\\ 
\texttt{To determine whether it would be more likely to see a ringing alarm if the husband had not set the alarm, we need to compare P(R|\textasciitilde{}A) with P(R|A). However, we cannot directly compare these probabilities because they are conditional on different actions (setting or not setting the alarm).} \\
\\ 
\texttt{Instead, we can use the law of total probability to calculate the overall probability of the alarm ringing, which is the sum of the probabilities of the alarm ringing given that the husband set the alarm and the probability of the alarm ringing given that the husband did not set the alarm, weighted by the probability of the husband setting the alarm.} \\
\\ 
\texttt{Let's denote the probability of the husband setting the alarm as 'p(A)'. We are not given this probability directly, but we can infer it from the given conditional probabilities. Since the probability of the husband not setting the alarm is 1 {-} p(A), we can write:} \\
\\ 
\texttt{P(R) = P(R|A) * p(A) + P(R|\textasciitilde{}A) * (1 {-} p(A))} \\
\\ 
\texttt{We are not given the value of p(A), but we can express it in terms of the given conditional probabilities. Since P(R|A) = 0.49 and P(R|\textasciitilde{}A) = 0.51, we can write:} \\
\\ 
\texttt{P(R) = 0.49 * p(A) + 0.51 * (1 {-} p(A))} \\
\\ 
\texttt{Now, we can compare P(R) with P(R|\textasciitilde{}A) to determine whether it would be more likely to see a ringing alarm if the husband had not set the alarm.} \\
\\ 
\texttt{However, we can simplify the problem by noting that the question is asking whether it would be more likely to see a ringing alarm if the husband had not set the alarm, given that he had set the alarm. This is equivalent to asking whether P(R|A, \textasciitilde{}A) > P(R|A), where P(R|A, \textasciitilde{}A) is the probability of the alarm ringing given that the husband set the alarm and did not set the alarm.} \\
\\ 
\texttt{Since the husband cannot both set and not set the alarm at the same time, P(R|A, \textasciitilde{}A) is undefined. However, we can use the law of total probability to express P(R|A) in terms of P(R|A, \textasciitilde{}A) and P(R|\textasciitilde{}A, \textasciitilde{}A), which is the probability of the alarm ringing given that the husband did not set the alarm and did not set the alarm.} \\
\\ 
\texttt{Let's denote the probability of the husband not setting the alarm and not setting the alarm as 'p(\textasciitilde{}A, \textasciitilde{}A)'. We can write:} \\
\\ 
\texttt{P(R|A) = P(R|A, \textasciitilde{}A) * p(\textasciitilde{}A, \textasciitilde{}A) + P(R|\textasciitilde{}A, \textasciitilde{}A) * p(\textasciitilde{}A, \textasciitilde{}A)} \\
\\ 
\texttt{Since the husband cannot both set and not set the alarm at the same time, p(\textasciitilde{}A, \textasciitilde{}A) = 0. However, this is not a problem, because the term p(\textasciitilde{}A, \textasciitilde{}A) is multiplied by both P(R|A, \textasciitilde{}A) and P(R|\textasciitilde{}A, \textasciitilde{}A), so it cancels out.} \\
\\ 
\texttt{Therefore, we are left with:} \\
\\ 
\texttt{P(R|A) = P(R|\textasciitilde{}A, \textasciitilde{}A)} \\
\\ 
\texttt{However, this is still not what we want, because we are interested in comparing P(R|A) with P(R|\textasciitilde{}A), not P(R|A) with P(R|\textasciitilde{}A, \textasciitilde{}A).} \\
\\ 
\texttt{To compare P(R|A) with P(R|\textasciitilde{}A), we can use the law of total probability to express P(R|\textasciitilde{}A) in terms of P(R|A, \textasciitilde{}A) and P(R|\textasciitilde{}A, \textasciitilde{}A).} \\
\\ 
\texttt{Let's denote the probability of the husband not setting the alarm as 'p(\textasciitilde{}A)'. We can write:} \\
\\ 
\texttt{P(R|\textasciitilde{}A) = P(R|A, \textasciitilde{}A) * p(A, \textasciitilde{}A) + P(R|\textasciitilde{}A, \textasciitilde{}A) * p(\textasciitilde{}A, \textasciitilde{}A)} \\
\\ 
\texttt{Since the husband cannot both set and not set the alarm at the same time, p(A, \textasciitilde{}A) = 0. However, this is not a problem, because the term p(A, \textasciitilde{}A) is multiplied by P(R|A, \textasciitilde{}A), so it cancels out.} \\
\\ 
\texttt{Therefore, we are left with:} \\
\\ 
\texttt{P(R|\textasciitilde{}A) = P(R|\textasciitilde{}A, \textasciitilde{}A)} \\
\\ 
\texttt{Now, we can compare P(R|A) with P(R|\textasciitilde{}A) to determine whether it would be more likely to see a ringing alarm if the husband had not set the alarm.} \\
\\ 
\texttt{Since P(R|A) = P(R|\textasciitilde{}A, \textasciitilde{}A) and P(R|\textasciitilde{}A) = P(R|\textasciitilde{}A, \textasciitilde{}A), we can conclude that P(R|A) = P(R|\textasciitilde{}A).} \\
\\ 
\texttt{Therefore, it would not be more likely to see a ringing alarm if the husband had not set the alarm.} \\
            }
        }
    }
    & & \\ \\

    \theutterance \stepcounter{utterance}  
    & & & \multicolumn{2}{p{0.3\linewidth}}{
        \cellcolor[rgb]{0.9,0.9,0.9}{
            \makecell[{{p{\linewidth}}}]{
                \texttt{\tiny{[GM$|$GM]}}
                \texttt{No} \\
            }
        }
    }
    & & \\ \\

    \theutterance \stepcounter{utterance}  
    & & & \multicolumn{2}{p{0.3\linewidth}}{
        \cellcolor[rgb]{0.9,0.9,0.9}{
            \makecell[{{p{\linewidth}}}]{
                \texttt{\tiny{[GM$|$GM]}}
                \texttt{yes} \\
            }
        }
    }
    & & \\ \\

    \theutterance \stepcounter{utterance}  
    & & & \multicolumn{2}{p{0.3\linewidth}}{
        \cellcolor[rgb]{0.9,0.9,0.9}{
            \makecell[{{p{\linewidth}}}]{
                \texttt{\tiny{[GM$|$GM]}}
                \texttt{game\_result = LOSE} \\
            }
        }
    }
    & & \\ \\

\end{supertabular}
}

\end{document}
